Conformal filters for vision--language model (VLM) captions keep an object claim only when an open-vocabulary detector finds enough evidence for it, and guarantee that at most a fraction $\delta$ of hallucinated claims is kept, provided test data are exchangeable with calibration data. Deployed systems break that assumption. We pre-registered and ran an audit of a claim-level and an image-level filter, both calibrated on clean COCO captions from LLaVA-1.5-7B scored by OWLv2, under 25 image degradations, two detector swaps, two VLM swaps and a shift to AMBER. All degradations preserved the claim-level target, although a one-sided stochastic-dominance condition that suffices for this held in only \NHtwo{} of them. Swapping the detector for Grounding DINO broke both filters (claim-level risk \GdinoClaim{} against a target of 0.10); recalibrating the threshold on the same labelled claims repaired them without new labels. Swapping the VLM broke the claim-level guarantee (InternVL3-2B: \InternClaim{}) but not the image-level one: the newer VLM hallucinates less often, but its hallucinations are more often objects the detector partly sees. A root-dominance bound for Ising trees under a shared perturbation budget explains why coupling a caption's claims in a graphical model cannot make the filters more robust at the coupling strength we measure.
